\section{Introduction}

\subsection{Policy context}

India operates one of the largest higher-education systems in the world, yet participation remains low relative to the size of its youth cohort. The National Education Policy 2020 sets a target of 50\% GER in higher education, including vocational education, by 2035 \cite{nep2020}. The official all-India GER was 28.4\% in 2021-22, an improvement of about one percentage point over 2020-21 but still roughly 22 points short of the goal \cite{aishe2023,ogd2025}. The Covid-19 period disrupted both enrolment and the administrative data series, so the recent trend is an uncertain guide to the next fourteen years.

The primary users of this analysis are State higher-education departments, which decide where to sanction new colleges and how to allocate faculty, and the central NEP monitoring cell, which tracks whether the 50\% target is reachable. The decision is not whether to expand---the target already commits the system to expansion---but where additional capacity would raise participation most per rupee. The question also sits within the international Sustainable Development Goal (SDG) 4 framework: target 4.3 commits governments to equal access for all women and men to affordable, quality technical, vocational and tertiary education by 2030 \cite{unesco2020}. According to the same UNESCO monitoring report, global tertiary GER reached 39\% (growing by roughly one percentage point per year since 2000), placing India's 28.4\% over ten points below the global benchmark, while participation in low- and middle-income countries remains heavily concentrated among the richest households \cite{unesco2020}. India's State-level gap is therefore both a national policy problem and part of the wider SDG 4 agenda.

The national figure is an average over units with very different starting points. In 2021-22, State GER ranged from 16.9\% in Assam to 47.0\% in Tamil Nadu; among Union Territories, Chandigarh reported 64.8\% and Dadra and Nagar Haveli and Daman and Diu 11.2\%. Fourteen of the 28 States were below 30\%, and none of the States had reached 50\%. Because education is on the concurrent list and because the demographic weight sits in the low-GER States, the national target is effectively a problem of State-level convergence.

\subsection{Research questions}

The project addresses three questions.

\begin{enumerate}
  \item \textbf{Association.} Conditional on the other indicators, how much independent information do college density, PTR and gender parity carry about State GER?
  \item \textbf{Stability.} Do those associations survive out-of-sample prediction, influence diagnostics and within-State panel variation, or are they artefacts of a 28-unit cross-section?
  \item \textbf{Projection.} Under explicit and auditable assumptions about density expansion, teacher supply and the common participation trend, which States reach 50\% by 2035-36, and what does that imply for the national target?
\end{enumerate}

The design is deliberately modest about causality. A cross-section of 28 States cannot identify the causal effect of building a college. It can, however, quantify conditional associations, test their fragility, and show how much the 2035 answer depends on assumptions that are usually left implicit.

\subsection{Contributions}

The study makes five methodological and reproducibility contributions.

\begin{itemize}
  \item It constructs an auditable, India-specific dataset from official government sources, cross-checked against published AISHE tables and stored with SHA-256 checksums and an explicit missingness audit.
  \item It implements a leakage-free nested cross-validation protocol across linear, regularised and non-parametric learners, and demonstrates why a better-fitting polynomial specification must be rejected for out-of-range extrapolation.
  \item It pairs point estimates with HC3 standard errors, 2,000-resample bootstrap intervals, variance inflation factors, residual specification tests, skewness/log-transform checks and leave-one-out influence analysis.
  \item It tests cross-sectional patterns against five-year within-State panel specifications and documents why the sign reversal on PTR serves as a cautionary negative finding.
  \item It provides a transparent, change-based projection framework with Monte Carlo uncertainty propagation and population-weighted aggregation to contrast State counts against cohort coverage.
\end{itemize}

Section~\ref{sec:related} reviews related work. Section~\ref{sec:data} describes the data and preprocessing, Section~\ref{sec:methods} the methods, Section~\ref{sec:setup} the experimental setup, Section~\ref{sec:results} the results, Section~\ref{sec:discussion} the discussion, and Section~\ref{sec:limitations} the limitations. Section~\ref{sec:conclusion} concludes.
